\chapter{Essential Algebra}
\paragraph{} Duckie and Güs's various competitions made the goose authorities a bit confused. Because of Duckie and Güs's competitions, they had lost track of Duckie's location. They decided that they needed to review their records, and try to figure out where their activities fifth day had landed them. 
\vfill
\pagebreak
%Equality
\subchapter{Equality}
{Duckie and Güs had tied in each of competitions. Because of this, their positions were always the same.}
{Güs's position = Duckie's position. Duckie's distance from the village was $\mathbf{24}km$, so Güs's position was also $\mathbf{24}km$.}
{When two things are the same, we write this in mathematics as "=", or "equals".}
{\begin{tikzpicture}
    \coordinate (A) at (0,0);
    \coordinate (B) at (0,6.4);
    \coordinate (Ap) at (2,0);
    \coordinate (Bp) at (2,6.4);

    \node [inner sep=0pt] at (B) {\includegraphics[height=0.8cm]{Gus}};
    \draw[->,ultra thick,red, decorate, decoration={random steps,segment length=3pt,amplitude=0.2pt}] (A) -- (B) node[midway, left] {$24 km$};
    \node [inner sep=0pt] at (Bp) {\includegraphics[height=0.8cm]{DuckieGami}};
    \draw[->,ultra thick,blue, decorate, decoration={random steps,segment length=3pt,amplitude=0.2pt}] (Ap) -- (Bp) node[midway, left] {$24 km$};
\end{tikzpicture}
}
\subchapter{The Equality Rule}
{The goose authorities, knew Güs was a fast flier. His training as a policegoose had made him a fearsome foe. While Duckie and Güs were traveled together, they realized that every time Güs flew ahead of Duckie, Duckie had to increase his position by the same amount to keep up.}
{Duckie's position + $\Delta$ position = Güs's position + $\Delta$ position.\footnote{$\Delta$ means "change in"}}
{What happens to one side of an equation must happen to the other side so that they are still the same.}
{\begin{tikzpicture}
    \coordinate (A) at (0,0);
    \coordinate (B) at (0,6.4);
    \coordinate (C) at (0,9.4);
    \coordinate (Ap) at (2,0);
    \coordinate (Bp) at (2,6.4);
    \coordinate (Cp) at (2,9.4);
    
    \node [inner sep=0pt] at (C) {\includegraphics[height=0.8cm]{Gus}};
    \draw[->,ultra thick,red, decorate, decoration={random steps,segment length=3pt,amplitude=0.2pt}] (A) -- (B) node[midway, left] {$24 km$};
    \draw[->,ultra thick,orange, decorate, decoration={random steps,segment length=3pt,amplitude=0.2pt}] (B) -- (C) node[midway, left] {$16 km$};
    \node [inner sep=0pt] at (Cp) {\includegraphics[height=0.8cm]{DuckieGami}};
    \draw[->,ultra thick,blue, decorate, decoration={random steps,segment length=3pt,amplitude=0.2pt}] (Ap) -- (Bp) node[midway, left] {$24 km$};
    \draw[->,ultra thick,orange, decorate, decoration={random steps,segment length=3pt,amplitude=0.2pt}] (Bp) -- (Cp) node[midway, left] {$16 km$};
    
    \draw [decorate, decoration = {calligraphic brace,mirror, raise=10pt,amplitude=5pt}] (Ap) -- (Cp) node[midway, right, xshift=0.5cm,line width=3pt] {$40 km$};
\end{tikzpicture}}
\subchapter{Variables}
{The goose authorities were still a bit confused about Duckie and Güs's progress. They knew that the two had $\mathbf{4}$ competitions, and that Duckie and Güs began competing when they were $\mathbf{24} km$ away from the village. They also knew that Duckie and Güs stopped competing at $\mathbf{40} km$. How much did they travel on average in each competition?}
{\begin{center}$4x + 24 = 40 \linebreak 4x + 24 = 40 - 24 \linebreak 4x = 16 \linebreak \frac{4x}{4} = \frac{16}{4} \linebreak x = 4  $\end{center}}
{A letters or symbol can be used to represent a number which isn’t already known, such as "$x$". $4x + 16 = 40$ means that $4\ast(some number) + 16 = 40$.}
{\begin{tikzpicture}
    \coordinate (A) at (0,0);
    \coordinate (B) at (0,1);
    \coordinate (C) at (0,2);
    \coordinate (D) at (0,3);
    \coordinate (E) at (0,4);
    \coordinate (F) at (0,10);
    \coordinate (Ap) at (1.5,0);
    \coordinate (Bp) at (1.5,1);
    \coordinate (Cp) at (1.5,2);
    \coordinate (Dp) at (1.5,3);
    \coordinate (Ep) at (1.5,4);
    \coordinate (Fp) at (1.5,10);
    \coordinate (Af) at (3.4,0);
    \coordinate (Bf) at (3.4,1);
    \node [inner sep=0pt] at (F) {\includegraphics[height=0.8cm]{DuckieGami}};
    \draw[->,ultra thick,blue, decorate, decoration={random steps,segment length=3pt,amplitude=0.2pt}] (A) -- (B) node[midway, left] {$? km$};
    \draw[->,ultra thick,blue, decorate, decoration={random steps,segment length=3pt,amplitude=0.2pt}] (B) -- (C) node[midway, left] {$? km$};
    \draw[->,ultra thick,blue, decorate, decoration={random steps,segment length=3pt,amplitude=0.2pt}] (C) -- (D) node[midway, left] {$? km$};
    \draw[->,ultra thick,blue, decorate, decoration={random steps,segment length=3pt,amplitude=0.2pt}] (D) -- (E) node[midway, left] {$? km$};
    \draw[->,ultra thick,green, decorate, decoration={random steps,segment length=3pt,amplitude=0.2pt}] (E) -- (F) node[midway, left] {$24 km$};
    \draw [decorate, decoration = {calligraphic brace, raise=35pt,amplitude=5pt}] (A) -- (F) node[midway, left, xshift=-1.5cm,line width=3pt] {$40 km$};
    
    \draw[->,ultra thick,blue, decorate, decoration={random steps,segment length=3pt,amplitude=0.2pt}] (Ap) -- (Bp) node[midway, left] {$? km$};
    \draw[->,ultra thick,blue, decorate, decoration={random steps,segment length=3pt,amplitude=0.2pt}] (Bp) -- (Cp) node[midway, left] {$? km$};
    \draw[->,ultra thick,blue, decorate, decoration={random steps,segment length=3pt,amplitude=0.2pt}] (Cp) -- (Dp) node[midway, left] {$? km$};
    \draw[->,ultra thick,blue, decorate, decoration={random steps,segment length=3pt,amplitude=0.2pt}] (Dp) -- (Ep) node[midway, left] {$? km$};
    \draw[->,ultra thick,red, decorate, decoration={random steps,segment length=3pt,amplitude=0.2pt}] (Fp) -- (Ep) node[midway, left] {$-24 km$};
    \draw [decorate, decoration = {calligraphic brace,mirror, raise=10pt,amplitude=5pt}] (Ap) -- (Ep) node[midway, right, xshift=0.5cm,line width=3pt] {$16 km$};
    
    \draw[->,ultra thick,blue, decorate, decoration={random steps,segment length=3pt,amplitude=0.2pt}] (Af) -- (Bf) node[midway, left] {$4 km$};
\end{tikzpicture}}.